% Einheit 3 von 4 – Einheit 3 (bank/flaecheninhalt-und-volumen-im-raum/e3.jsonl, zusammenbau v0.1)
%% TODO Prüfungswort Einheit 3: nicht in der Marken-Zeile
%% TODO Zeitmarke Einheit 3: Marken-Zeile nicht lesbar
%% TODO Zweigzeile Teil 1: was hier gelernt wird (Satz in Schülersprache aus dem Einheitstitel) – Einheit 3
%% TODO Einheitentitel 3 nicht in der Mappe lesbar
\einheitenkopf[e3]{Einheit 3 von 4 · Einheit 3}
\setcounter{aufgabe}{14}

%% TODO Titel: Ich-kann-Satz fehlt in der Bank (Platzhalter: Kettenname) – Nr. 15
%% TODO Anweisung: ein Satz über der ersten Teilaufgabe fehlt in der Bank – Nr. 15
\begin{aufgabe}{Volumen und Oberfläche}
%% TODO \rechenplatz{n}: Schrittzahl des Grundfalls fehlt in der Bank (nur bei fester Schreibform, 2.3 b)
\begin{teile}
\teil Ein Zelt über einem quadratischen Boden, alle Dachkanten laufen in einer Spitze zusammen. Volumenformel? \\ \kreuz{mit Faktor ein Drittel} \\ \kreuz{ohne Faktor ein Drittel}
\teil Pyramide mit der Grundfläche $P(1 | 1 | 0)$, $Q(5 | 1 | 0)$, $R(5 | 5 | 0)$, $T(1 | 5 | 0)$ und der Spitze $S(3 | 3 | 6)$. Volumen? V = \leerfeld VE
\teil Das Dreieck $OPQ$ mit $P(0 | 3 | 4)$ und $Q(6 | 0 | 0)$ ist rechtwinklig in $O$. Die Pyramide $OPQS$ hat das Volumen 40, und es gilt $\overrightarrow{OS} \cdot \overrightarrow{OP} = 0$ und $\overrightarrow{OS} \cdot \overrightarrow{OQ} = 0$. Länge $|OS|$? \hfill (Abitur 2026 GK) \\ |OS| = \leerfeld
\teil Die Pyramide $PQRTS$ hat das Drachenviereck mit $P(0 | 0 | 0)$, $Q(3 | 2 | 0)$, $R(0 | 7 | 0)$, $T(-3 | 2 | 0)$ als Grundfläche und die Spitze $S(0 | 0 | 5)$. Volumen? \hfill (Abitur 2025 GK) \\ V = \leerfeld VE
\teil $P(1 | 1 | 0)$, $Q(5 | 1 | 0)$, $R(1 | 4 | 0)$, $P'(1 | 1 | 6)$, $Q'(5 | 1 | 6)$, $R'(1 | 4 | 6)$. Weise nach, dass $PQRP'Q'R'$ ein gerades Prisma ist, und berechne sein Volumen. V = \leerfeld VE
\teil Pyramide mit der quadratischen Grundfläche $P(-3 | -3 | 1)$, $Q(5 | -3 | 1)$, $R(5 | 5 | 1)$, $T(-3 | 5 | 1)$ und der Spitze $S(1 | 1 | 4)$. Oberflächeninhalt? \hfill (Abitur 2024 LK) \\ O = \leerfeld FE
\teil $P(1 | 0 | 2)$, $Q(4 | 4 | 2)$, $R(-4 | 10 | 2)$, $T(-7 | 6 | 2)$, $S(-1 | 5 | 8)$. Zeige, dass $PQRT$ ein Rechteck, aber kein Quadrat ist, und berechne das Volumen der Pyramide $PQRTS$.
\end{teile}
\end{aufgabe}

%% TODO Titel: Ich-kann-Satz fehlt in der Bank (Platzhalter: Kettenname) – Nr. 16
%% TODO Anweisung: ein Satz über der ersten Teilaufgabe fehlt in der Bank – Nr. 16
\begin{aufgabe}{Volumen und Oberfläche – Fehler finden, Begründen, Anwendung}
\begin{teile}
\teil Eine Pyramide hat die Grundfläche 24 und die Höhe 5. Lena rechnet: \rechnung{V &= 24 \cdot 5 \\ &= 120} Finde den Fehler und rechne richtig.
\teil Begründe: Eine Pyramide hat ein Drittel des Volumens eines Prismas mit gleicher Grundfläche und gleicher Höhe.
\teil Das Obergeschoss eines Veranstaltungssaals ist die Pyramide mit der waagerechten dreieckigen Grundfläche $P(0 | 0 | 12)$, $Q(0 | 24 | 12)$, $R(-20 | 6 | 12)$ und der Spitze $S(-8 | 8 | 30)$ (1 LE = 1 m). Die Lüftung braucht 0{,}5\,kW je 100\,m³. Reicht eine Anlage mit 8\,kW? \hfill (Abitur 2018 LK)
\end{teile}
\end{aufgabe}
